\documentclass{library/duthesis}

\documenttype{Thesis}

\title{A Step towards the Origin of Fracture}
\author{Samuel Walker}
\dept{Department of Physics}
\degree{Doctor of Philosophy}
\degreedate{March 2026}
\copyrightyear{2026}

%% File to be included while running latex.         
%\includeonly{chapters/1-Introduction/chapter,       
%             chapters/2-Background/chapter,          
%             references/ref,                        
%             appendices/a/appendix}                 

\begin{document} 
% Frontmatter.
\frontmatter
\titlepage
\newpage

\mainmatter
\chapter{Introduction} 


\chapter{Background}
\section{Rheology}
\begin{itemize}
    \item What is rheology
    \item Definitions of stress and strain from an intuitive/geometric definition ($\epsilon = \frac{\Delta x}{L}$ and $\sigma = \frac{F}{A}$)
\end{itemize}
\subsection{Strain}
\begin{itemize}
    \item Derivation of the Cauchy’s Finite strain tensor from first principles in terms of the deformation gradient tensor $\bm{F}$
    \item Discussion on linear elasticity theory to get $\epsilon_{ij}$
    \item Define Elogational Strain for Context and Simple (single and EPM) and Pure shear (Double Network) for Results chapters - Use Diagrams here
\end{itemize}
\subsection{Stress}
\begin{itemize}
    \item Define Stress tensor (Which definition/derivation to use?).
    \item Go on define Generalised Hooks law and apply isotropic and incompressible constraints
\end{itemize}

\section{Hydrodynamics}
\begin{itemize}
    \item Navier Stokes $\rightarrow$ Stokes equation $\rightarrow$ force balance
    \item Splitting of strain field into elastic and \quotes{plastic} fields as in Picard et al
    \item Derivation of Oseen/Stokeslet 
    \item Construction of the Elshelby Stress tensor as in Picard et al
    \item Mention of streamline approximation
\end{itemize}
\section{Networks}
\subsection{Semi-Flexible filaments}
\begin{itemize}
    \item Discussion of Semi-Flexible filaments
    \item Discussion of Hamiltonian from MacKintosh et al
    \item Take a central force model, ignoring the bending force
\end{itemize}
\subsection{Rigidity}
\begin{itemize}
    \item While not directly explored in the thesis, the concept underpins a lot of the physics/choices made. 
    \item Define rigidity
    \item Disscuss maxwell criteria and introduce average connectivity $z$
    \item Where does Maxwell's criteria not work - Over constrained systems (regions of self stress) and Non-homogeneous networks
    \item Relation to percolation
    \item Types of rigidity
    \begin{itemize}
        \item First Order rigidity - from network structure
        \item Second Order rigidity - Arises from force balances and non-linear effects
    \end{itemize}
    \item Energetic Rigidity - thinking of rigidity in terms of energy landscapes
\end{itemize}
\subsection{Fracture}
\begin{itemize}
    \item Discussion on how fracture is handled in network models
\end{itemize}
\subsection{Discorded Network topologies}
\begin{itemize}
    \item Lattice-based networks - So I can add a little bit of the early work on double networks
    \item Jamming-based networks - Description of the generation of a single network, as this covers the fundamentals of both network chapters
    \item Dangling and Disconnected clusters and how to prune networks
\end{itemize}
\subsection{Viral Stress tensor}
\begin{itemize}
    \item Derivation of the Viral Stress tensor
\end{itemize}
\subsection{Lees Edwards Boundary conditions}
\begin{itemize}
    \item How shear is applied in MD simulation through Lees-Edwards boundary conditions - Mostly maths with maybe a bit of pseudo code - might include triclinic boundary sets, which is the mathematical equivalent
\end{itemize}

\chapter{Fatigue in Lattice EPM}
\section{Introduction}
\begin{itemize}
    \item Discussion of LAOS Protocol and why it is uses in experiments
    \item Discussion of previous results in Oscillatory shear
    \item What is an absorbed state, and what does it mean to yield in a quasistatic EPM model
\end{itemize}

\section{Methodology}
\begin{itemize}
    \item Implementation of a quasistatic Athermal EPM using results in the introduction
\end{itemize}

\section{Results}
\begin{itemize}
     \item Discussion on what an Absorbed state vs a Yielded state looks like ,with snapshots of the systems
    \item Distributions of what state the systems end up in using stacked bar charts
    \item Survivalship PDF Plots showing the probabilities of reaching the delayed dividing, Couple this with a slice of each when less than 10\% of samples pass this point showing???? behaviour
    \item Long delayed yielding is driven by a decay to a random walk of one oscillator, where we have a two thresholds as to when the system absorbs or Yields
    \item Maybe Hyper-uniform or distributions here, depending on results
\end{itemize}

\chapter{Thermal Delayed Yielding in Networks}
\section{Introduction}
\begin{itemize}
    \item Follows work done previously in continuum models by Lockwood and Carrington
    \item Discussion on what the expected stress response of the system is
    \begin{itemize}
        \item Would expect either a shear band to form immediately as the material is taken above the critical point
        \item Or a slow relaxation to the equilibrium state
        \item However, it is possible that the failure can be delayed almost indefinitely, with thermal activity eventually causing the failure
        \item Highly delayed banding under step strain has been less explored.
        \item This is important as the system exists in an apparent steady state (equilibrium point) for a very long time before thermal energy causes failure, so it could be unexpected if no knowledge of the strain history is known
    \end{itemize}
    \item Experimental work
    \item Discussion on shear Banding vs Fracture
    \item We are going to use two models to explore this one thermal model and an Athermal model 
\end{itemize}

\section{Brownian Model}
\begin{itemize}
    \item Change to model (add noise term)
    \item Reference Thermal Stiffening - See Tauber et al 2020
    \item Effect of $\gamma_0$ and temperature $T$
    \item Computational cost required due to a small fixed time step means we can't access ultra-long time scales.
\end{itemize}

\section{SGR-Like Breakage Condition}
\begin{itemize}
    \item Implementation of Breakage condition in SGR-like way, Equivalent to a Boltzmann distribution
    \item Analogy of the particles in a bond sitting in a potential well, randomly exploring until it exits and the bond breaks - A Boltzmann distribution drives the rate at which this happens.
    \item Show early time relaxation to the same state after an imposed strain
    \item Effect of Temperature on breakage time $t^*$, time is delayed with an exponential factor, amount of broken bonds converges to a constant small amount as the fracture behaviour goes from diffuse to localised
    \item effect of $\gamma_0$ on the fracture time as we get closer to the failure point, The cracks again go from a diffuse to a localised regime
    \item Direction of crack propagation changes as a function of $\gamma_0$. In the small strain regime, the system cracks evenly vertically or horizontally with even probability. At a large strain, it cracks diagonally, following the principal compression direction perpendicular to the axis of shear.
\end{itemize}

\section{Discussion}
\begin{itemize}
    \item Discussion on Athermal vs Thermal, Brownian is a thermal simulation, Breakage Condition added an effective temperature to an Athermal model
\end{itemize}


\chapter{Critical Bonds in Networks}
\section{Introduction}
\begin{itemize}
    \item In the cases where we get a single crack at failure, is it possible to predict before failure where the crack might form
    \item It would be great to understand any precursors to failure which has wide-reaching real-world applications: a good understanding of the precursors to yielding could allow engineers to identify when a material is likely to fail, even if there is no obvious indication from its macroscopic stress-strain relationship.
\end{itemize}
\section{Results}
\begin{itemize}
    \item Network Strain Distributions as a function of $\gamma_0$
    \item Break Angles of Bonds in Simulations
    \item Size of Avalanches after artificially breaking a bond
    \item Edge Betweenness centrality
\end{itemize}



\chapter{Toughness of double network hydrogels: the role of reduced
stress propagation}
\section{Introduction}
\begin{itemize}
    \item What is a hydrogel? - context to polymer networks
    \item Challenge: polymer networks are typically either brittle and stiff or ductile and flexible
    \item Examples of why we might want a stiff/tough double network material
    \begin{itemize}
        \item tendons, ligaments and cartilage
        \item Scaffolds for tissue engineering
        \item Load-bearing gels
    \end{itemize}
    \item Explain pre-existing experiential work by gong, creton, ect. 
    \item Lack of existing theoretical work on Double networks
    \item Aim is to create a simple mesoscopic model
\end{itemize}

\section{Methodology}
\begin{itemize}
    \item Extension of the single network model to the double network
    \begin{itemize}
        \item Square lattice - choice of overlapping bonds
        \item Jammed Network - Overlapping on nodes
    \end{itemize}
    \item Pure shear protocol - How the Quasistatic regime is implemented
    \item Averaging of stress tensor to obtain heat maps
\end{itemize}

\section{Results}
\subsection{Square Networks}
\begin{itemize}
        \item Stress Strain Curves in Square Network
        \item effect of \quotes{overlapping} bonds
        \item Propagator in Square network
\end{itemize}
\subsection{Jammed Networks Networks}
\begin{itemize}
    \item Affect of $\mu$ and $M$ on stress strain curve responses
    \item Compromise this lead to between $\epsilon^*$, $\Sigma^*$ and $G_s/G_d$
    \item Diffuse vs Single crack failure as we add the second network
    \item Propagator in Jammed network and the reduction in the propagator as a function of $\mu$ and $M$
    \item Discuss the effect of $z$ on the model behaviour.
\end{itemize}
\end{document}

